\documentclass{article}

\usepackage[margin=1in]{geometry}
\usepackage[linesnumbered,ruled,vlined]{algorithm2e}
\usepackage{amsmath,amssymb,amsthm}
\usepackage{stmaryrd} % for double brackets
\usepackage{mathrsfs} % good script letters

% wrap in angle brackets
\newcommand{\tuple}[1]{\langle #1 \rangle}

% interpretation of regular expression as a transition formula
\newcommand{\tf}[1]{\mathscr{I} \llbracket #1 \rrbracket}
\newcommand{\si}[1]{\mathscr{L} \llbracket #1 \rrbracket}

% (strictly)dominates
\newcommand{\strictlydominates}{\mathrel{\gg}}
\newcommand{\dominates}{\mathrel{\underline{\strictlydominates}}}
\newcommand{\strictlydominatedby}{\mathrel{\ll}}
\newcommand{\dominatedby}{\mathrel{\underline{\strictlydominatedby}}}


% define theorems, lemmas, corollaries, and definitions
\newtheorem{theorem}{Theorem}
\newtheorem{lemma}{Lemma}
\newtheorem{corollary}{Corollary}

\theoremstyle{definition}
\newtheorem{definition}{Definition}

\begin{document}

\begin{definition}
  The set of \emph{regular expressions} over $\Sigma$, written
  $RegExp(\Sigma)$, is the smallest set such that
  $0, 1 \in RexExp(\Sigma)$, $\Sigma \subseteq RegExp(\Sigma)$,
  and $(R_1 + R_2), (R_1 \cdot R_2), R^* \in RegExp(\Sigma)$ for
  all $R, R_1, R_2 \in RegExp(\Sigma)$.

  $SubExp(R)$ for $R \in RegExp(\Sigma)$ is the set of
  subexpressions of $R$, including $R$ itself.
\end{definition}

\begin{definition}
  % taken more or less verbatim from Algebraic Program Analysis +
  % Fast Algorithms for Solving Path Problems
  The \emph{standard interpretation} is the function
  $\si{\cdot}: RegExp(\Sigma) \to \textbf{L}$, where $\textbf{L}$
  is the set of regular languages, such that
  $\si{0} = \emptyset$, $\si{1} = \{\epsilon\}$ where $\epsilon$
  is the empty string, $\si{a} = \{a\}$ for all $a \in \Sigma$,
  and
  \begin{align*}
  \si{R_1 + R_2} &= \si{R_1} \cup \si{R_2} \\
  \si{R_1 \cdot R_2} &= \{x_1x_2| x_1 \in \si{R_1}, x_2 \in
                       \si{R_2}\} \\
  \si{R^*} &= \{x_1x_2\dots x_n| x_1,x_2,\dots,x_n \in \si{R}\}
  \end{align*}
  for all $R, R_1, R_2 \in RegExp(\Sigma)$.
\end{definition}

\begin{definition}
  Given a function $wt: \Sigma \to \mathbb{Z}$, the \emph{length}
  of a regular expression is given by
  $L: RegExp(\Sigma) \to \mathbb{Z} \cup \{-\infty\}$ where
  $L(0) = -\infty$, $L(1) = L(R^*) = 1$, $L(a) = wt(a)$ for all
  $a \in \Sigma$, $L(R_1 + R_2) = \max\{L(R_1), L(R_2)\}$, and
  $L(R_1 \cdot R_2) = L(R_1) + L(R_2)$.
\end{definition}

Let $X$ and $X'$ be sets of variables such that for each
$x\in X$, there is a corresponding $x' \in X'$. As an abuse of
notation, we use $X = X' \equiv \bigwedge_{x \in X} x =
x'$. $Vars(F)$ for a formula $F$ is the set of free variables in
$F$. $Formula(X, X')$ denotes the set of formulas such that
$Vars(F) \subseteq X \cup X'$. 

\begin{definition}
  A \emph{control flow graph} $G = \tuple{V, E, r, \mathscr{F}}$
  is a graph with vertices $V$, edges $E$, a root vertex $r$ of
  indegree 0 from which every vertex $v \in V$ is reachable, and
  a function $\mathscr{F}: E \to Formula(X, X')$ for some
  $X, X'$. We use $\mathscr{F}_{\tuple{u,v}}(Y, Y')$ to denote
  $\mathscr{F}(\tuple{u,v})[X \mapsto Y, X' \mapsto Y']$. We let
  $Paths_G(u,v)$ be the set of all (possibly non-simple) paths
  from $u$ to $v$ in $G$.
\end{definition}

\begin{definition}
  A \emph{path expression} $P$ of type $\tuple{u, v}$ on
  $G = \tuple{V, E, r, F}$ is a regular expression over
  $E$ such that $\si{P} \subseteq Paths_G(u,v)$ for some
  $u,v \in V$. We denote the set of path expressions on $G$ by
  $PathExp(G)$.
\end{definition}

Suppose $star$ is a function that takes a formula $F$ and the
pair of variable vocabulairies $X$ and $X'$ that $F$ relates and
returns some overapproximation of the reflexive transitive
closure of $F$ over $X$ and $X'$ containing no other free
variables.

\begin{definition}
  The \emph{transition formula interpretation} is the function
  $\tf{\cdot}: PathExp(G) \to Formula(X, X')$ where
  $\tf{0} \equiv false$, $\tf{1} \equiv X = X'$,
  $\tf{\tuple{u, v}} \equiv \mathscr{F}_{\tuple{u, v}}(X, X')$,
  and
  \begin{align*}
    \tf{P_a + P_b} &\equiv \tf{P_a} \lor \tf{P_b} \\
    \tf{P_a \cdot P_b} &\equiv \exists X''. \tf{P_a}(X, X'')
                         \land \tf{P_b}(X'', X') \\
    \tf{P^*} &\equiv star(\tf{P}, \tuple{X, X'}) 
  \end{align*}
  for all $P, P_a, P_b \in PathExp(G)$. 

  We use $\tf{P}(Y, Y')$ to denote
  $\tf{P}[X \mapsto Y, X' \mapsto Y']$.
\end{definition}

\begin{definition}
  Suppose there is an infinite set of variable vocabularies
  $\{X_0, X_1, X_2, \dots\}$, and suppose that $L$ is taken with
  respect to the function $wt: E \to \mathbb{Z}$ that maps every
  edge to 1. Then, define $I$ mapping path expressions and
  integer lower and upper bounds to formulas where
  $F(P, l, u) \equiv false$ if $L(P) = -\infty$, and otherwise,
  $F(1, l, u) \equiv X_l = X_u$,
  $F({\tuple{u, v}}, l, u) \equiv \mathscr{F}_{\tuple{u, v}}(X_l,
  X_u)$, and
  \begin{align*}
    F(P_a + P_b, l, u)
    &\equiv F(P_a, l, u) \lor F(P_b, l, u) \\
    F(P_a \cdot P_b, l, u)
    &\equiv F(P_a, l, l + L(P_a)) \land F(P_b, l + L(P_a), u) \\
    F(P^*, l, u)
    &\equiv star(F(P, l+1, l+1+L(P)), \tuple{X_{l+1},
      X_{l+1+L(P)}})[X_{l+1} \mapsto X_l, X_{l+1+L(P)} \mapsto X_u]
  \end{align*}
  for all $P, P_a, P_b \in PathExp(G)$. 
\end{definition}

\begin{lemma}
  \label{pe:freevars}
  For all $P \in PathExp(G)$, for all $l \ge 0$ and
  $u \ge l + L(P)$,
  $Vars(F(P, l, u)) \subseteq \bigcup_{i = l}^u X_i$.
\end{lemma}

\begin{proof}
  The claim holds trivially when $L(P) = -\infty$, $P = 1$,
  $P = P'^*$, or $P = \tuple{v, w}$. Suppose the claim holds for
  $P_a$ and $P_b$ and that $L(P_a + P_b) \ge 0$. Take any
  $0 \le l \le l + L(P_a + P_b) \le u$. Since
  $L(P_a + P_b) = \max\{L(P_a), L(P_b)\}$, clearly
  $u \ge l + L(P_a)$ and $u \ge l + L(P_b)$. Then,
  $Vars(F(P_a + P_b, l, u)) \subseteq Vars(F(P_a, l, u)) \cup
  Vars(F(P_b, l, u)) = \bigcup_{i = l}^u X_i$. Similarly, suppose
  $L(P_a \cdot P_b) \ge 0$. Then, $L(P_a), L(P_b) \ge 0$. Take
  $0 \le l \le l + L(P_a \cdot P_b) = 1 + L(P_a) + L(P_b) \le u$.
  Clearly, $u \ge (l + L(P_a)) + P_b$. So,
  $Vars(F(P_a \cdot P_b, l, u)) \subseteq Vars(F(P_a, l, l +
  L(P_a))) \cup Vars(F(P_b, l + L(P_a), u)) = (\bigcup_{i = l}^{l
    + L(P_a)} X_i) \cup (\bigcup_{i = l + L(P_a)}^u X_i) =
  \bigcup_{i = l}^u X_i$. By induction, the claim must hold for
  all $P$.
\end{proof}

\begin{theorem}
  For all $P \in PathExp(G)$,
  $Sat(\tf{P}) \iff Sat(F(P, 0, L(P)))$.
\end{theorem}

\begin{proof}
  Let $V_{i,j} = \bigcup_{k=i}^j X_k$. As an abuse of notation,
  let $\exists S$ for a set $S = \{x_1, \dots, x_m\}$ denote
  $\exists x_1, \dots, x_m$. Then, we proceed by structural
  induction on $P$ with the hypothesis that for all $l \ge 0$ and
  $u \ge l + L(P)$,
  $\tf{P} \equiv \exists V_{l,u}. F(P, l, u) \land X = X_l \land
  X' = X_r$. By Lemma \ref{pe:freevars}, $X$ and $X'$ are not
  free in $F(P, l, u)$, so this is sufficient to prove the claim.
  If $L(P) = -\infty$, $P = 1$, or $P = \tuple{v, w}$, then the
  claim holds trivially.

  Suppose the claim holds for $P_a$ and $P_b$ and suppose
  $L(P_a + P_b) \ge 0$. Take arbitrary
  $0 \le l \le l + L(P_a + P_b) \le u$. Then,
  \begin{align*}
    \tf{P_a + P_b}
    &\equiv \tf{P_a} \lor \tf{P_b} \\
    &\equiv
      (\exists V_{l,u}. F(P_a, l, u) \land X = X_l \land X' = X_r)
      \lor (\exists V_{l,u}. F(P_b, l, u) \land X = X_l \land X' =
      X_r) \\
    &\equiv
      \exists V_{l,u}. F(P_a+P_b,l,u) \land X = X_l \land X' = X_r
  \end{align*}

  Similarly, suppose $L(P_a \cdot P_b) \ge 0$ and take arbitrary
  $0 \le l \le l + L(P_a \cdot P_b) \le u$. By Lemma
  \ref{pe:freevars}, we know that
  $Vars(F(P_a, l, l + L(P_a))) \subseteq V_{l, l + L(P_a)}$ and
  $Vars(F(P_b, l + L(P_a), u)) \subseteq V_{l + L(P_a), u}$.  By
  definition,
  $V_{l, u} = V_{l, l + L(P_a)} \cup V_{l + L(P_a), u}$ and
  $V_{l, l + L(P_a)} \cap V_{l + L(P_a), u} = X_{l +
    L(P_a)}$. Furthermore, $L(P_a), L(P_b) \ge 0$ by definition,
  meaning that $l + L(P_a) \ge 0$ and
  $u \ge (l + L(P_a)) + L(P_b)$. So,
  \begin{align*}
    \tf{P_a \cdot P_b}
    &\equiv \exists X''.
      \tf{P_a}(X, X'') \land \tf{P_b}(X'', X') \\
    &\equiv \exists X''.
      (\exists V_{l, l + L(P_a)}. F(P_a, l, l + L(P_a)) \land
      X = l \land X'' = X_{l + L(P_a)}) \\
    & \qquad \land (\exists V_{l + L(P_a), u}.
      F(P_b, l + L(P_a), u) \land
      X'' = X_{l + L(P_a)} \land X' = u) \\
    &\equiv \exists X''.
      \exists V_{l, u}.
      F(P_a, l, l + L(P_a)) \land F(P_b, l + L(P_a), u) \land
      X = l \land X'' = X_{l + L(P_a)} \land X' = u \\
    &\equiv \exists V_{l, u}.
      F(P_a \cdot P_b, l, u) \land X = l \land X' = u
  \end{align*}

  Finally, suppose the claim holds for $P$ and take arbitrary
  $0 \le l \le l + 1 \le u$. By the induction hypothesis,
  $\tf{P} \equiv \exists V_{l+1, l+1+L(P)}. F(P, l+1, l+1+L(P))
  \land X = X_{l+1} \land X' = X_{l+1+L(P)}$. So,
  $\tf{P^*}(X_{l+1}, X_{l+1+L(P)}) \equiv star(F(P, l+1, l+1+L(P)),
  \tuple{X_{l+1}, X_{l+1+L(P)}})$. Therefore,
  $\tf{P^*}(X_l, X_u) \equiv F(P^*, l, u)$, and the claim holds.
  
\end{proof}

\begin{definition}
  In a control flow graph $G = \tuple{V, E, r, \mathscr{F}}$,
  $v\in V$ \emph{dominates} $u \in V$, written $v \dominates u$,
  if every path from $r$ to $u$ in $G$ goes through $v$. If
  $v \dominates u$ and $u \ne v$, then $v$ \emph{strictly
    dominates} $u$, written $v \strictlydominates u$. If $v$
  strictly dominates $u$ and is dominated by every vertex that
  strictly dominates $u$, then $v$ is the \emph{immediate
    dominator} of $u$. Note that every vertex besides $r$ has a
  unique immediate dominator --- we denote this by
  $idom(u)$. $E|_v$ is the set of edges $\tuple{u,w}$ such that
  $v \strictlydominates w$. $\tuple{u, v} \in E$ is a \emph{back
    edge} if $v \dominates u$.
\end{definition}

\begin{lemma}
  \label{dom:trans}
  If $u \dominates v \dominates w$, then every path
  from $u$ to $w$ must go through $v$.
\end{lemma}

\begin{proof}
  If $u = v$ or $v = w$, then the claim holds trivially.
  Otherwise, let $p_1$ be a simple path from $r$ to $u$ and let
  $p_2$ be an arbitrary path from $u$ to $w$. By definition, the
  concatenation of $p_1$ and $p_2$ must go throught $v$. Suppose
  $p_1$ goes through $v$. Then, there is some prefix $p_1'$ of
  $p_1$ that is a path from $r$ to $v$. But, by definition,
  $p_1'$ must go through $u$, which contradicts the simplicity of
  $p_1$. Therefore, $p_2$ must go through $v$, as desired.
\end{proof}

\begin{lemma}
  \label{dom:nostrict}
  If $u \not\strictlydominates v$, $v \not\strictlydominates u$,
  $u \not\strictlydominates w$, and $v \dominates w$, then every
  path from $u$ to $w$ must pass through $v$.
\end{lemma}

\begin{proof}
  If $u = v$ or $v = w$, then the claim holds trivially. If
  $u = w$, then $v = w$, so the claim also holds
  trivially. Otherwise, consider a path $p_1$ from $r$ to $u$
  that doesn't pass through $v$. Then, take any path $p_2$ from
  $u$ to $w$. The concatenation of $p_1$ and $p_2$ is a path from
  $r$ to $w$, so it must pass through $v$. So, $p_2$ passes
  through $v$.
\end{proof}

\begin{definition}
  A control flow graph $G = \tuple{V, E, r, \mathscr{F}}$ is
  \emph{reducible} if for every cycle $v_0,v_1,\dots,v_n=v_0$ in
  $G$, there is some $0 \le i \le n$ such that
  $v_i \dominates v_j$ for all $0 \le j \le n$. (Intuitively,
  every cycle has a natural entrypoint.)
\end{definition}

\begin{lemma}
  \label{reducible:order}
  If $G$ is reducible, then there is an ordering of the vertices
  such that for all $u, v\in V$, $u \le v$ if there is a vertex
  $u'$ where $u \dominatedby u' \not \strictlydominates v$ and
  there is a path from $u'$ to $v$ that doesn't traverse any back
  edges.
\end{lemma}

\begin{proof}
  We prove the claim by showing that on reducible graphs, the
  above relation is a partial order.
  \begin{enumerate}
  \item Reflexivity: let $v \in V$ be arbitrary.
    $v \dominatedby v \not \strictlydominates v$, and the empty
    path from $v$ to $v$ does not traverse any back edges.
  \item Antisymmetry: suppose there are some $u, u', v, v' \in V$
    and paths $p_1$ and $p_2$ between $u'$ and $v$ and $v'$ and
    $u$, respectively, such that
    $u \dominatedby u' \not \strictlydominates v$,
    $v \dominatedby v' \not \strictlydominates u$, and $p_1$ and
    $p_2$ don't traverse any back edges. Since
    $u' \not\strictlydominates v$ and
    $v' \not\strictlydominates u$, neither $u'$ nor $v'$ can
    strictly dominate the other. So, by Lemma \ref{dom:nostrict},
    $p_1$ must go through $v'$, and $p_2$ must go through
    $u'$. But, this means there is a cycle starting and ending at
    $u'$ that goes through $v'$ and doesn't traverse a back
    edge. If the cycle is non-empty, then no vertex dominates its
    predecessor in the cycle. Since $G$ is reducible, this means
    the cycle must be empty, i.e. $u' = v'$. So,
    $u' \dominates v$ but $u' \not \strictlydominates v$ and
    $v' \dominates u$ but $v' \not \strictlydominates
    u$. Therefore, $u = v' = u' = v$.
  \item Transitivity: suppose there are some
    $u, u', v, v', w \in V$ and paths $p_1$ and $p_2$ between
    $u'$ and $v$ and $v'$ and $w$, respectively, such that
    $u \dominatedby u' \not \strictlydominates v$,
    $v \dominatedby v' \not \strictlydominates w$, and $p_1$ and
    $p_2$ don't traverse any back edges. Note that
    $u' \not \strictlydominates v'$. If $v' \dominates u'$, then
    we are done. Otherwise, $v' \not \strictlydominates u'$. By
    Lemma \ref{dom:nostrict}, $p_1$ passes through $v'$. So,
    there is a path from $u'$ to $w$ through $v'$ that doesn't
    traverse any back edges. But, this also means that
    $u' \not \strictlydominates w$, because otherwise
    $u' \dominates v'$.
  \end{enumerate}
\end{proof}

\begin{algorithm}[H]
  \caption{Tarjan's Path Expression Algorithm on Reducible Flow Graphs}
  \label{alg:tarjan}

  \SetAlgoLined
  % \SetAlgoNoEnd
  
  \KwIn{A reducible control flow graph
    $G=\tuple{V, E, r, \mathscr{F}}$ with vertices labeled
    $1,\dots, n = r$ as detailed in Lemma \ref{reducible:order}}
  
  \KwOut{A function $pe: V \to PathExp(G)$ such that for all
    $v \in V$, $pe(v)$ is a well-formed path expression
    recognizing exactly $Paths_G(r, v)$}

  Initialize a compressed weighted forest on $V$ \;

  \For{$v = 1$ \KwTo $n-1$}{
    $cycle, noncycle \gets 0$ \;

    \ForEach{$v' \in pred(v)$}{
      \lIf{$find(v') = v$}{
        $cycle \gets cycle + eval(v') \cdot \tuple{v', v}$
      }
      \lElse{
        $noncycle \gets noncycle + eval(v') \cdot \tuple{v', v}$
      }
    }

    $link(v, noncycle \cdot cycle^*, idom(v))$ \;
  }

  \Return $\lambda v. eval(v)$ \;

\end{algorithm}

Note that constant time simplication procedures applied after
each concatenation or union can easily remove all 0's and 1's
from the path expressions --- we assume that this is the case for
the rest of the analysis of Tarjan's algorithm.

\begin{lemma}
  \label{tarjan:correct}
  Before the $v$th iteration, for all $w < v$, $find(w) = D_v(w)$
  and $\si{eval(w)} = Paths_G(D_v(w), w) \cap E|_{D_v(w)}^*$,
  where $D_v(w)$ is the smallest vertex at least $v$ that
  dominates $w$.
\end{lemma}

\begin{proof}
  We proceed by induction on $v$. The claim holds vacuously
  before the first iteration. Suppose it holds before the $v$th
  iteration. By the induction hypothesis, $find(w) = v$ only if
  $v \dominates w$. So, for all $w \not \dominatedby v$,
  $find(w)$ and $eval(w)$ are unchanged, and
  $D_{v+1}(w) = D_v(w)$. So, the claim still holds for such
  $w$. Consider any $w \dominatedby v$. By the ordering,
  $w \le v$, and $D_v(w) = v$. By the induction hypothesis,
  before this iteration,
  $\si{eval(w)} = Paths_G(v, w) \cap E|_v^*$. Furthermore,
  $D_{v+1}(w) = idom(v)$, and every path from $idom(v)$ to $w$
  must pass through $v$ by Lemma \ref{dom:trans}. Thus, it is
  sufficient to show that
  $\si{noncycle \cdot cycle^*} = Paths_G(idom(v), v) \cap
  E|_{idom(v)}^*$. That the left hand side is a subset of the
  right hand side follows trivially from the induction
  hypothesis, so we focus on the other direction.

  Consider any such path $p$. It can be chopped into paths
  $p_0, p_1,\dots,p_m$ where no $p_i$ for $0 \le i \le m$
  contains $v$ as an intermediate vertex. $p_0$ can be further
  expressed as the concatenation of a path $p_0'$ and an edge
  $\tuple{v', v}$ where $v' \dominatedby idom(v)$. By Lemma
  \ref{dom:trans}, since $p_0'$ doesn't go through $v$,
  $v \not \dominates v'$. And, since $v' \dominatedby idom(v)$,
  $v'$ must be dominated by every vertex that strictly dominates
  $v$. So, either $v' = idom(v)$, or
  $v' \not \strictlydominates v$ and $\tuple{v', v}$ is a path
  from $v'$ to $v$ that doesn't traverse a back edge, i.e.
  $v' \le v$. Suppose there is some $v''$ such that
  $v' \dominatedby v'' \not \dominates v$. There must be a path
  from $v''$ to $v'$ that does not traverse a back edge. So,
  $v'' < v$. Therefore, the smallest vertex at least $v$ that
  dominates $v'$ is $idom(v)$. By the induction hypothesis,
  $p_0'$ is recognized by $eval(v')$. So, $noncycle$ recognizes
  $p_0$.

  Now, take any of the $p_i$. By reducibility, some vertex in
  $p_i$ must dominate all the others. In particular, it must
  dominate $v$. But, $p_i$ does not pass through
  $idom(v)$. Therefore, $v$ must dominate all vertices in
  $p_i$. In particular, we can express $p_i$ as the concatenation
  of a path $p_i'$ and an edge $\tuple{v', v}$, where
  $v' \dominatedby v$. By the induction hypothesis, $p_i'$ is
  recognized by $eval(v')$ before this iteration. So, $cycle$
  recognizes $p_i$ for all $1 \le i \le m$. Therefore, the claim
  holds after this iteration.
\end{proof}

\begin{definition}
  If $P$ is a path expression such that
  $0, 1 \not \in SubExp(P)$, then let $src, tgt: SubExp(P) \to V$
  and $sstar, tstar: SubExp(P) \to \{true, false\}$ be defined as
  follows:

  \begin{center}
    \begin{tabular}{ c | c | c | c | c }
      & $src$ & $tgt$ & $sstar$ & $tstar$ \\
      \hline
      $\tuple{v, w}$ & v & w & $false$ & $false$ \\
      $P_1 + P_2$ & $src(P_1)$ & $tgt(P_1)$ & $sstar(P_1) \lor sstar(P_2)$ & $tstar(P_1) \lor tstar(P_2)$ \\
      $P_1 \cdot P_2$ & $src(P_1)$ & $tgt(P_2)$ & $sstar(P_1) $ & $tstar(P_2)$ \\
      ${P_0}^*$ & $src(P_0)$ & $tgt(P_0)$ & $true$ & $true$
    \end{tabular}
  \end{center}

  We let
  $type(P_0) = \tuple{\tuple{src(P_0), sstar(P_0)},
    \tuple{tgt(P_0), tstar(P_0)}}$ and
  $starred(P_0) \equiv sstar(P_0) \land tstar(P_0)$.
\end{definition}

\begin{lemma}
  Suppose $eval_v(\cdot)$ is the state of the compressed weighted
  forest before the $v$th iteration. Let
  $S_v = \bigcup_{w < v} SubExp(eval_v(w))$,
  $K_v^t = \max\{L(P)|P \in S_v, type(P) = t\}$, and
  \begin{align*}
    l_{v,k}(P)
    &= \begin{cases}
      k+L(eval_v(src(P))) & \neg starred(P) \\
      k+L(eval_v(tgt(P))) - 1 & starred(P)
    \end{cases} \\
    u_{v,k}(P) &= l_{v,k}(P) + K_v^{type(P)}
  \end{align*}
  for any $k \ge 0$. Then, for all $k \ge 0$, the following hold:

  \begin{enumerate}
  \item For all $P = P_1 + P_2 \in S_v$,
    $F(P, l_{v,k}(P), u_{v,k}(P)) \equiv F(P_1, l_{v,k}(P_1),
    u_{v,k}(P_1)) \lor F(P_2, l_{v,k}(P_2), u_{v,k}(P_2))$.
  \item For all $P = P_1 \cdot P_2 \in S_v$,
    $F(P, l_{v,k}(P), u_{v,k}(P)) \equiv F(P_1, l_{v,k}(P_1),
    u_{v,k}(P_1)) \land F(P_2, l_{v,k}(P_2), u_{v,k}(P_2))$
  \item For all $P = {P_0}^* \in S_v$ with $l = l_{v,k}(P)$,
    $u = u_{v,k}(P)$, $l_0 = l_{v,k}(P_0)$, and
    $u_0 = u_{v,k}(P_0)$,
    $F(P, l, u) \equiv star(F(P_0, l_0, u_0), \tuple{X_{l_0},
      X_{u_0}})[X_{l_0} \mapsto X_l, X_{u_0} \mapsto X_u]$
  \end{enumerate}
\end{lemma}

\begin{proof}
  We proceed by induction on $v$. The claim holds vacuously for
  $v = 0$. Suppose it holds for $0 \le v \le n - 1$ and take an
  arbitrary $k \ge 0$. By Lemma \ref{tarjan:correct},
  $eval_{v+1}(w) = eval_v(w)$ if $v \not \dominates w$, and
  $eval_{v+1}(w) = noncycle \cdot cycle^* \cdot eval_v(w)$
  otherwise. Let
  $S_1 = \bigcup_{w < v, w\not\dominatedby v} SubExp(eval_v(w))$,
  $S_2 = \bigcup_{w < v, w\strictlydominatedby v}
  SubExp(eval_v(w))$, and $S_3 = SubExp(noncycle \cdot
  cycle^*)$. Then, $S_v = S_1 \cup S_2$ and
  $S_{v+1} = S_v \cup S_3$.

  Take any $P_1 \in S_1$. There must be some
  $w \not \dominatedby v$ such that $P_1 \in
  SubExp(eval_v(w))$. By Lemma \ref{tarjan:correct},
  $w \dominates init(P_1), fin(P_1)$. Then,
  $v \not \dominates init(P_1), fin(P_1)$. So,
  $l_{v+1,k}(P_1) = l_{v,k}(P_1)$ and
  $u_{v+1,k}(P_1) = u_{v,k}(P_1)$. Thus, the claim holds for $S_1$.

  Similarly, take any $P_2 \in S_2$. By Lemma
  \ref{tarjan:correct}, $v \dominates init(P_2), fin(P_2)$. Let
  $L = L(noncycle \cdot cycle^*)$. Then,
  $L(eval_{v+1}(init(P_2))) = L + L(eval_v(init(P_2)))$ and
  $L(eval_{v+1}(fin(P_2))) = L + L(eval_v(fin(P_2)))$. So,
  $l_{v+1,k}(P_2) = L + l_{v,k}(P_2) = l_{v,k+L}(P_2)$ and
  $u_{v+1,k}(P_2) = L + u_{v,k}(P_2) =
  u_{v,k+L}(P_2)$. Therefore, the claim holds for $S_2$.

  It remains to show that the claim holds for all $P_3 \in
  S_3$. For all $v' \in pred(v)$ where $v'\ne idom(v)$,
  $sstar(eval(v')) = tstar(\tuple{v',v}) = false$. Furthermore,
  $eval(idom(v))\cdot \tuple{idom(v), v}$ gets simplified to
  $\tuple{idom(v), v}$, and
  $sstar(\tuple{idom(v), v}) = tstar(\tuple{idom(v), v}) =
  false$. Therefore, $K_v^{type(noncycle)} = L(noncycle)$ and
  $K_v^{type(cycle)} = L(cycle)$. So, (1) holds for all $S_3$;
  (2) holds for $noncycle \cdot cycle^*$ since
  $L(eval_{v+1}(v)) = L(noncycle) + 1$ and $starred(cycle^*)$;
  and (3) holds for $cycle^*$ since $starred(cycle^*)$ but
  $\neg starred(cycle)$.

  Thus, we only need to show that (2) holds for all
  $eval_v(v') \cdot \tuple{v', v}$ where $v' \in pred(v)$. Let
  $v'$ be arbitrary and take the smallest $w'$ at least $v$ that
  dominates $v'$. Since $v' \ne v$, the only $P_3\in S_v$ with
  $src(P_3) = w'$ and $tgt(P_3) = v'$ is $eval_v(v')$. So,
  $K_v^{type(eval_v(v'))} = L(eval_v(v'))$. Therefore, (2) holds
  for $eval_v(v') \cdot \tuple{v', v}$, as desired.
\end{proof}

\end{document}
